\SourceLeaf{16}
\par\medskip\Needspace{4\baselineskip}
\centerline{THEOREM SIX.}\medskip
\par
{\itshape The attraction that the body exerts on point $M$ in the direction parallel and opposite to the $x$-coordinate axis is expressed by the integral
\[
-\int\frac{\mathrm{d}s.\Phi r.\cos MQ.\cos MX}{rr}
\]
taken over the whole surface, whether $M$ lies outside the body or inside it. The term $-\pi\Phi0.\cos QX$ must be added whenever $M$ lies on the surface itself, in which case the particular value that $QX$ has at $M$ is to be used.}\par
Evidently the forces in the directions parallel and opposite to the coordinate axes $y,z$ are likewise expressed by the integrals
\[
-\int\frac{\mathrm{d}s.\Phi r.\cos MQ.\cos MY}{rr},\quad-\int\frac{\mathrm{d}s.\Phi r.\cos MQ.\cos MZ}{rr}
\]
to which $-\pi\Phi0.\cos QY$, $-\pi\Phi0.\cos QZ$ must be added (taking the particular values of the angles at $M$) whenever $M$ lies on the surface of the body.

It is also easy to see that the three forces
\[
-\pi\Phi0.\cos QX,\quad-\pi\Phi0.\cos QY,\quad-\pi\Phi0.\cos QZ
\]
are equivalent to a single force $=-\pi\Phi0$ normal to the surface itself and directed inward.

{\rightskip=0pt plus 3em We could evidently have dispensed with evaluating the integral $-\Phi0.\int\mathrm{d}\Sigma.\cos MX$ if the function $f$ had the appropriate form. It would have to allow us to set $\Phi0=0$. But we preferred to pursue the investigation in full generality. Whenever the attraction is assumed to be inversely proportional to the cube or a higher power of the distance, it is clear that this is impossible and that we necessarily have $\Phi0=-\infty$. It follows that under this assumption a point on the surface of the body is pressed toward the solid with infinite force.\par}

\Article{9}
\par
The integrals we have considered would have had to be taken over the whole volume of the body (triple integrals). By the methods explained so far, we have reduced them, in two different ways, to integrals that need only be taken over the surface of the body. The form of the surface is expressed by an equation between the coordinates $x,y,z$, that is, by the equation $W=0$, where $W$ denotes a function of the variables $x,y,z$ that may be assumed free of irrational terms. Let the differentiation of the function $W$ give
\[
\mathrm{d}W=T\mathrm{d}x+U\mathrm{d}y+V\mathrm{d}z
\]
